% Results-section addendum: counterfactual experiments.
% Drafted 2026-05-13 for inclusion at the end of Section~\ref{sec:results},
% after the hukou subsubsection (sec:hukou) and after the interpretation
% addendum (sec:interpretation, see results_interpretation_addendum.tex).
% Standalone .tex; paste directly into sec_results.tex.
%
% Cross-references used:
%   eq:decision-rule, eq:restricted-GRC, eq:gmm-moments, eq:phi-proposition1
%     (defined in sec_model.tex; verified).
%   sec:hukou, sec:grc-returns, sec:interpretation (sec:results.tex).
%   app:inversion-preview (sec_model/appendix; verified).
% Citations used (verified against CKT.bib):
%   bryanAggregateProductivityEffects2019, lagakosMigrationCostsObservational2020,
%   lagakosUrbanRuralGapsDeveloping2020a, suriSelectionComparativeAdvantage2011.

\subsection{Counterfactual experiments}\label{sec:counterfactuals}

The restricted GRC delivers a worker-specific return $\Delta_i = \beta + \phi\theta_i$ and a full set of trajectory-conditional averages $\Delta_{\underline{d}}$.
We use these objects to address two questions that translate the slope and intercept of our LCA line into magnitudes the migration literature presses on.
First, how much aggregate consumption is left on the table by the gap between observed sorting and the sorting that an optimizing planner with knowledge of $\Delta_i$ would implement?
Second, how much of that gap in China is attributable to hukou-induced restrictions on rural-urban mobility, and what would removing those restrictions deliver?
Both experiments stay on the consumption side of $\Delta_i$, which is what the LCA framework identifies cleanly; we return to the welfare interpretation at the end of the section.

Two features of this exercise are worth flagging up front.
The exercise sits in a tradition of structural migration counterfactuals---\cite{bryanAggregateProductivityEffects2019} on the aggregate productivity gain from removing rural-urban frictions in Indonesia, \cite{lagakosMigrationCostsObservational2020} on the migration costs implied by observed wage gaps in panel data, and \cite{lagakosUrbanRuralGapsDeveloping2020a} on the urban-rural productivity gap more broadly---but with a specific feature that our identification delivers: trajectory-conditional returns identified from switchers and extrapolated to non-switchers through the LCA restriction, with inference propagated through the inversion confidence region rather than asymptotic standard errors that the M\"obius singularity at $\phi = -1$ would make unreliable.
The counterfactuals also inherit the testable piece of our identification: in specifications where the $J$-test rejects, the magnitudes below are reported with the same caution we apply to the underlying $\Delta_{d_T}$ and $\Delta_{d_N}$ estimates, and we report results separately by hukou regime where the rejection is regime-driven.

\subsubsection{Aggregate consumption gap from misallocated migration}\label{sec:misallocation}

The aggregate misallocation gap brackets observed migration against two extremes.
At one end, a world with no rural-urban migration, where workers stay in the location of their initial trajectory observation regardless of $\Delta_i$.
At the other end, a world of optimal sort, where each worker spends all periods in the location that maximizes their consumption ($D_{it} = 1$ iff $\Delta_i > 0$, $D_{it} = 0$ otherwise).
Observed migration sits between the two.
The contrast between the three scenarios is the readable summary statistic the counterfactual delivers.

Let $W_{\text{obs}}, W_{\text{zero}}, W_{\text{opt}}$ denote per-capita log consumption in the observed, zero-migration, and optimal-sort scenarios.
With trajectory population shares $\pi_{\underline{d}}$, urban time shares within trajectory $\bar{D}_{\underline{d}}$, and trajectory-conditional returns $\Delta_{\underline{d}}$, the two differences are
\begin{equation}\label{eq:misallocation-decomposition}
    \underbrace{W_{\text{obs}} - W_{\text{zero}}}_{\text{value of observed migration}} = \sum_{\underline{d}} \pi_{\underline{d}} \, \Delta_{\underline{d}} \, \bar{D}_{\underline{d}}, \qquad
    \underbrace{W_{\text{opt}} - W_{\text{obs}}}_{\text{misallocation gap}} = \sum_{\underline{d}} \pi_{\underline{d}} \left[\max(0, \Delta_{\underline{d}}) - \Delta_{\underline{d}} \, \bar{D}_{\underline{d}}\right].
\end{equation}
The first piece values the migration that actually occurs.
The second piece measures what is left on the table because workers with positive returns failed to move, or because workers with negative returns moved.
Reporting both lets readers locate observed migration on the available frontier: a small ``value of observed migration'' alongside a large ``misallocation gap'' would sharpen our pro-poor finding into a policy conclusion that the workers who would gain most from migration are precisely those who do not migrate.

Each term in equation~\eqref{eq:misallocation-decomposition} is identified by the GRC.
The trajectory shares $\pi_{\underline{d}}$ come from the descriptive panel.
The urban time shares $\bar{D}_{\underline{d}}$ are deterministic functions of trajectory definitions.
The switcher returns $\Delta_{\underline{d}}$ for $\underline{d} \in \mathcal{D}_S$ are identified non-parametrically from the unrestricted GRC.
The always-urban return $\Delta_{d_T}$ is identified from the LCA inversion, and the never-migrant return $\Delta_{d_N} = \beta + \phi(\mu_{d_N} - \mu_{\underline{d}_0})$ is identified by the LCA extrapolation in equation~\eqref{eq:restricted-GRC}.
Inference for the aggregate proceeds by recomputing equation~\eqref{eq:misallocation-decomposition} at each $(\phi, \beta)$ in the joint inversion confidence region (Appendix~\ref{app:inversion-preview}) and reporting the convex hull of the resulting gap distribution.
This propagates the uncertainty in $\phi$ and $\beta$ through the aggregate without requiring delta-method approximations that the M\"obius pole at $\phi = -1$ would not support.

A subtlety we will address explicitly is that the LCA restriction identifies trajectory-conditional means $\Delta_{\underline{d}}$ but not the within-trajectory dispersion of $\Delta_i$.
The trajectory-mean formula in equation~\eqref{eq:misallocation-decomposition} treats every never-migrant as if they share the return $\Delta_{d_N}$, which is conservative whenever $\Delta_{d_N} > 0$: by Jensen's inequality applied to $\max(0, \cdot)$, the within-trajectory expectation of the positive part weakly exceeds the positive part of the conditional mean.
We report three magnitudes that bracket what LCA can say about the gap.
First, a conservative floor that uses only trajectory means.
Second, a heterogeneity-corrected estimate that integrates $\max(0, \beta + \phi\theta)$ against the density of $\theta_i \mid d_N$ implied by the model's decision rule in equation~\eqref{eq:decision-rule}, under a stated parametric form for $\nu_{it}^l$.
Third, a bounded envelope that parameterizes within-trajectory dispersion as a fraction of the cross-trajectory dispersion and reports the gap as a function of that fraction.
The choice between the three is a choice about how much within-trajectory structure to assume; presenting all three lets the audience read the magnitude robust to that assumption.

Two design choices follow our reading of the existing literature.
We report the misallocation gap in log points and as a percentage of observed aggregate consumption, following \cite{bryanAggregateProductivityEffects2019}, who report counterfactual gains as percentage changes in productivity.
We also report a decomposition of the gap into the contribution of never-migrants ($d_N$), switchers, and always-urban workers ($d_T$), because under our pro-poor finding the never-migrant piece dominates the aggregate and the decomposition is the cleanest way to communicate where the misallocation lives.

\subsubsection{Removing the hukou wedge in China}\label{sec:hukou-counterfactual}

The hukou-split estimates in Section~\ref{sec:hukou} reveal a structural contrast that motivates a second counterfactual.
The rural-hukou subsample carries a large unrealized return on a flat LCA slope.
The urban-hukou subsample carries an essentially zero return on a steep slope.
Both regimes pass the $J$-test within their own subsample, so the LCA restriction is not refuted in either regime.
Read against the model's decision rule in equation~\eqref{eq:decision-rule}, the two estimates tell a coherent story: where the institutional barrier is absent, workers sort on comparative advantage and the residual never-migrants have approximately zero pecuniary returns to migration; where the barrier binds, sorting is suppressed and high-return rural-hukou workers remain rural.
The counterfactual quantifies what removing the barrier would deliver to the rural-hukou population.

We report two versions, parallel to the literature's bound-and-magnitude convention.
First, a lower bound from the identified objects alone.
Take the flat $\phi^{rh}$ at face value as genuine uniformity of returns within the rural-hukou regime, so that every rural-hukou never-migrant carries the same return $\Delta_{d_N}^{rh}$.
The aggregate consumption gain from moving the rural-hukou never-migrants to their optimal location is then
\begin{equation}\label{eq:hukou-bound}
    \text{Consumption gain (lower bound)} = \pi^{rh} \cdot \pi_{d_N}^{rh} \cdot \Delta_{d_N}^{rh},
\end{equation}
where $\pi^{rh}$ is the rural-hukou share of the CHN sample and $\pi_{d_N}^{rh}$ is the never-migrant share within the rural-hukou subsample.
The expression is a bound rather than a magnitude: if the flat $\phi^{rh}$ reflects suppressed sorting rather than uniformity of returns, the true gain from removing the barrier exceeds this floor because optimal sort would select the high-return tail of the rural-hukou never-migrant distribution rather than averaging across it.

Second, a resorting version that delivers a magnitude rather than a bound.
The model's decision rule in equation~\eqref{eq:decision-rule} already specifies how workers choose location given their type; we use it as written.
We parameterize the hukou barrier as an additive cost that shifts the rural-hukou intercept $\beta^{rh}$ down to the urban-hukou intercept $\beta^{uh}$, so removing the barrier sets the wedge to zero and the rule operates without institutional distortion.
The resorting version then simulates which rural-hukou workers would select urban under the no-barrier rule, computes the aggregate consumption gain, and decomposes that gain into the piece realized by the marginal migrant (who moves under no-hukou but did not move under hukou) and the piece left unrealized in the residual never-migrant pool (who stay rural even without the barrier because their personal $\Delta_i$ is close to zero or negative).
Inference propagates the inversion confidence intervals on $\phi^{rh}$ and $\phi^{uh}$ through the simulation.

Two open implementation choices will be reported transparently.
The rural-hukou and urban-hukou regimes are estimated with separate base trajectories $\underline{d}_0^{rh}$ and $\underline{d}_0^{uh}$, so the magnitude of the bound depends on the choice of base.
We will report results both with the original regime-specific bases used in Tables~\ref{tab:GRC_CHN_hukou_rural_first_consumption_urban_unb} and \ref{tab:GRC_CHN_hukou_urban_first_consumption_urban_unb} and with a common base re-estimated for the counterfactual, and report the range.
The resorting version also requires a choice for the distribution of $\nu_{it}^U - \nu_{it}^R$ beyond i.i.d.\ across individuals, time, and locations, which the model leaves open in Section~\ref{sec:model_identification}.
We report results under both a type-I extreme value distribution (logit-style choice probabilities, the most common choice in the empirical migration literature) and a normal distribution (probit-style, closer to the Heckman selection convention).
Reporting both gives a defensible range and lets the reader see how much the magnitude depends on the distributional assumption rather than the identification.

\subsubsection{From consumption to welfare}\label{sec:welfare-bridge}

Both experiments above report consumption-side statistics.
Converting them to welfare statements requires a value for the non-pecuniary component of $\nu_{it}^U - \nu_{it}^R$, which the LCA framework deliberately does not identify on its own.
A reallocation that raises consumption may also forgo non-pecuniary value that the consumption side alone cannot price, and the interpretation discussion in Section~\ref{sec:interpretation} qualifies the welfare reading accordingly.
Our reported counterfactuals should be read as the consumption-side complement to welfare statements: the headline number is what reallocation would deliver in consumption units, with the welfare gain weakly smaller whenever non-pecuniary value at the worker's observed location is positive.

A natural extension recovers the non-pecuniary side from observable correlates of location choice, following \cite{suriSelectionComparativeAdvantage2011}'s Section 7 strategy of relating residual choice probabilities to amenity proxies (distance to roads, distance to fertilizer sellers, infrastructure) rather than imputing them from a distributional assumption.
The analogous variables in our three panels---distance from origin village, family network at origin, ethnic and community ties, marital status---are available in CFPS, IFLS, and TZNPS, with country-specific augmentations (hukou in CHN, transmigration history in IDN, linguistic region in TZA).
The exercise generalizes the hukou comparison cross-country and pins down what fraction of the gap between observed sorting and rational-sorting predictions is attributable to identifiable amenity wedges versus residual unexplained variation.
We pursue this welfare extension in ongoing work.
